\section{Finite matrix kernels and positive trace}
\label{sec:weil-formal-kernels}

A scalar obtained by closing an operator kernel need not define a
positive pairing.  This distinction can already be seen on finite
matrices inside the formal operator algebra.
For this comparison, use the underlying field $\C$ with the discrete
topology and put
\[
 H_{\mathrm{for}}=\C[[x,y]]/\C.
\]
Enumerate the nonconstant monomials as $e_0,e_1,\ldots$ in increasing
total degree.  The formal topology makes $H_{\mathrm{for}}$ the product
$\C^{\mathbb N}$ of discrete coefficient spaces.  Write
$\epsilon^j(h)=h_j$ for its coefficient functionals.
A continuous endomorphism has a matrix with finitely many nonzero entries
in each row.  Indeed, every output coefficient must depend on finitely
many input coefficients.  Conversely, this condition makes every finite
output projection continuous.  Denote this algebra by $\mathscr E$.
Its operator topology requires agreement of each entire row in a fixed
finite set of output rows.

Put $E_{ij}=e_i\otimes\epsilon^j$ and let $\mathscr F_0$ be their
finite complex span.  These are the matrices with finitely many nonzero
entries in total.  They satisfy
\[
 E_{ij}E_{k\ell}=\delta_{jk}E_{i\ell},\qquad
 E_{ij}^*=E_{ji}.
\]
Conjugating coefficients defines the displayed involution on
$\mathscr F_0$.  This algebra has no unit, since each of its elements
annihilates all but finitely many basis vectors.
It is dense in $\mathscr E$: retaining finitely many whole rows of a
row-finite matrix gives a matrix in $\mathscr F_0$ with those rows unchanged.

The same matrix units act on $D_0=\C^{(\mathbb N)}$, the finite sequences,
inside the Hilbert space $\ell^2(\mathbb N)$ over the usual complex field.
Its inner product is $\langle v,w\rangle=\sum_jv_j\overline{w_j}$.
Every finite matrix preserves $D_0$, and conjugate transpose is its
Hilbert adjoint restricted to $D_0$.  Thus this comparison preserves the
algebra and involution on $\mathscr F_0$.  It does not identify the two
topologies or their completions.

\begin{proposition}[The positive finite-matrix trace]
\label{prop:weil-matrix-trace}
Define $\tau(B)=\sum_i B_{ii}$ on $\mathscr F_0$.
It is a positive linear functional with zero GNS null space.
Its GNS completion is $\ell^2(\mathbb N\times\mathbb N)$, with $[B]$
identified with its coefficient array $(B_{ij})$.
The GNS action is left matrix multiplication, and
\[
 \|\pi_\tau(A)\|=\|A\|_{\ell^2\to\ell^2}
 \qquad(A\in\mathscr F_0).
\]
Nevertheless, $\tau$ has no positive extension of finite mass to
$\mathscr F_0^+=\C I+\mathscr F_0$.
\end{proposition}
\begin{proof}
All sums in the algebraic calculation are finite, and
\[
 \tau(C^*B)=\sum_{i,j}B_{ij}\overline{C_{ij}},\qquad
 \tau(B^*B)=\sum_{i,j}|B_{ij}|^2.
\]
The second expression vanishes exactly when $B=0$.
Finite arrays are dense in the space of square-summable arrays by
truncation of their convergent sum of squared absolute values.
This identifies the quotient and completion.
Multiplication by $A$ acts separately on each column, so
\[
 \|AB\|^2=\sum_j\|A B_{\bullet j}\|_{\ell^2}^2
       \leq\|A\|_{\ell^2\to\ell^2}^2\sum_j\|B_{\bullet j}\|_{\ell^2}^2.
\]
The operator norm on the right is finite.  The entrywise
Cauchy--Schwarz inequality bounds it by
$(\sum_{i,j}|A_{ij}|^2)^{1/2}$.
Conversely, take a matrix $B$ whose only nonzero column is a finite
vector $v$.  The ratio $\|AB\|/\|B\|$ is $\|Av\|/\|v\|$.
Since $A$ uses only finitely many input coordinates, such vectors
compute its full operator norm.  This proves equality.

For $n\geq1$, let $Q_n=\sum_{i=0}^{n-1}E_{ii}$.
It is a self-adjoint projection with $\tau(Q_n)=n$.
If a positive extension had value $c$ on $I$, then $c$ would be real
and nonnegative, while
\[
 0\leq\tau^+((I-Q_n)^*(I-Q_n))=c-n
\]
for every $n$.  No finite $c$ satisfies these inequalities.
Equivalently,
$|\tau(Q_n)|^2/\tau(Q_n^*Q_n)=n$ is unbounded.
\end{proof}

The map $(\lambda,B)\mapsto\lambda I+B$ identifies the algebraic
unitization with the displayed subalgebra of $\mathscr E$.
It is injective because a scalar multiple of the infinite identity is
a finite matrix only when that scalar is zero.
On this subalgebra, any linear extension of finite-matrix trace has the
form
\[
 \tau_c(\lambda I+B)=c\lambda+\operatorname{tr}(B).
\]
It is cyclic, as finite matrix multiplication shows, but the proposition
excludes positivity for every $c$.  In particular, any cyclic trace on
a larger formal differential-operator algebra that restricts to this
formula cannot be positive for an involution restricting to the one above.
This conclusion needs no choice of an involution on the larger algebra.

The finite-dimensional corners $Q_n\mathscr F_0Q_n$ do have positive
traces, and the traces agree under their inclusions.  Their units are
$Q_n$, with masses $n$.  Compatibility of these finite traces supplies
$\tau$ on their union, but supplies no uniform mass bound on a unitization.
The common GNS domain and a vector representing the functional are
therefore separate constructions even in this explicit matrix example.

\begin{proposition}[Formal completion and Hilbert adjoints]
\label{prop:weil-formal-adjoint}
Conjugate transpose on $\mathscr F_0$ does not extend to a continuous
involution on $\mathscr E$ with its operator topology.
The finite-matrix action on $D_0$ also does not extend by the same
coefficient formula to an action of all $\mathscr E$ on
$\ell^2(\mathbb N)$.
\end{proposition}
\begin{proof}
Let $T_n=\sum_{i=0}^nE_{i0}$.
Every fixed row stabilizes, so $T_n$ converges in $\mathscr E$ to
\[
 T(h)=h_0(1,1,1,\ldots).
\]
The matrices $T_n^*=\sum_{i=0}^nE_{0i}$ do not converge in
$\mathscr E$.  Their zeroth rows do not stabilize, and an entire output
row is discrete in the operator topology.  Continuity of a proposed
involution would send the convergent sequence $T_n$ to a convergent
sequence, which is impossible.
For the second assertion, $e_0\in\ell^2$ but
$T(e_0)=(1,1,\ldots)\notin\ell^2$.
\end{proof}

There is also a direct norm obstruction.
The projections $E_{nn}$ tend to zero in the formal operator topology,
but each has Hilbert operator norm one and GNS vector norm one for
$\tau$.  Thus the formal completion, operator-norm completion and
GNS completion are different constructions on the same finite matrices.
For the Weil algebra, Theorem~\ref{thm:graph-domain} specifies a fourth
topology when its multipliers are unbounded: the common graph topology.
None of these completions is selected by cyclicity alone.

\section{Gaussian averages and formal coefficients}
\label{sec:weil-gaussian-coefficients}

The Gaussian in the explicit formula also permits an exact comparison
between a numerical integral and formal polynomial elimination.
It is necessary to specify both the observable algebra and the parameter.
For a positive real number $h$, set
\[
 \mathbb E_h(p)=\frac1{\sqrt{2\pi h}}
       \int_{\R}p(u)e^{-u^2/(2h)}\,du
       =\int_{\R}p(u)g_{h/2}(u)\,du,
 \qquad p\in\C[u].
\]
The polynomial algebra here uses pointwise multiplication and
$p^*(u)=\overline{p(u)}$ for real $u$.
Its involution fixes $u$.  Polynomial growth and Gaussian decay make
every displayed integral absolutely convergent.

\begin{proposition}[Numerical and formal Gaussian elimination]
\label{prop:weil-gaussian-comparison}
For every polynomial $p$ and real $h>0$,
\begin{equation}
 \mathbb E_h(p)
 =\sum_{m\geq0}\frac{h^m}{2^m m!}\,p^{(2m)}(0).
 \label{eq:weil-gaussian-moments}
\end{equation}
The sum is finite.  In particular, $\mathbb E_h$ is positive,
$\mathbb E_h(u)=0$ and $\mathbb E_h(u^2)=h$.

Let $\hbar$ instead be a formal central variable of degree zero.
The same finite derivative formula on each polynomial gives a continuous
$\C[[\hbar]]$-linear map
\[
 \mathsf G:\C[u][[\hbar]]\longrightarrow\C[[\hbar]],\qquad
 \mathsf G\!\left(\sum_{j\geq0}p_j(u)\hbar^j\right)
  =\sum_{j,m\geq0}\frac{p_j^{(2m)}(0)}{2^m m!}\hbar^{j+m}.
\]
Continuity refers to the $\hbar$-adic topologies.
On $\C[u,\hbar]$, evaluating this result at $\hbar=h>0$ agrees
with the numerical Gaussian average of the specialized polynomial.
There is no unital $\C$-algebra homomorphism
$\C[[\hbar]]\to\C$ that sends $\hbar$ to a nonzero number $h$.
\end{proposition}
\begin{proof}
Lemma~\ref{lem:gaussian} gives $\mathbb E_h(1)=1$.
Reflection makes every odd moment zero.
For $n\geq2$, integrate the derivative of
$u^{n-1}e^{-u^2/(2h)}$ over $\R$.
Its values at both ends are zero, and its derivative is integrable.
It follows that
\[
 \mathbb E_h(u^n)=(n-1)h\,\mathbb E_h(u^{n-2}).
\]
Thus
$\mathbb E_h(u^{2m})=(2m)!h^m/(2^m m!)$.
Expansion of $p$ proves \eqref{eq:weil-gaussian-moments}.
The positive integral of $|p|^2$ proves positivity, and the moments
of $u$ and $u^2$ have the stated values.

In the formal formula, the coefficient of $\hbar^N$ uses only
$j,m\geq0$ with $j+m=N$.  There are finitely many such pairs.
Each derivative is a derivative of a polynomial, so this defines every
coefficient without a convergence assumption.
The map preserves every power of $\hbar$, proving continuity.
Reindexing these finite coefficient sums proves
$\C[[\hbar]]$-linearity.  For a polynomial in $u,\hbar$, all sums are
finite, so numerical specialization commutes with the proved moment formula.

For $h\ne0$, the formal series
\[
 (1-\hbar/h)^{-1}=\sum_{n\geq0}(\hbar/h)^n
\]
is an inverse in $\C[[\hbar]]$.
A unital homomorphism sending $\hbar$ to $h$ would send the unit
$1-\hbar/h$ to zero, which is impossible.
At $h=0$, the constant-coefficient homomorphism does exist.
\end{proof}

The comparison at positive $h$ therefore holds on the polynomial
subalgebra.  Extension to a selected class of infinite series requires
an additional summability rule and its compatibility with the operations.
Formal coefficientwise existence alone gives no such rule.

The Gaussian average is not multiplicative:
$\mathbb E_h(u)^2=0$ while $\mathbb E_h(u^2)=h$.
Nor does inserting a Gaussian turn pointwise polynomial multiplication
into convolution.  For $A>0$, the linear map
\[
 J_A:\C[u]\longrightarrow\A_\zeta,\qquad J_A(p)(u)=p(u)g_A(u)
\]
is well-defined, since every derivative is a polynomial times a Gaussian.
But
\[
 J_A(1)\star J_A(1)=g_{2A}\ne g_A=J_A(1),
 \qquad J_A(u)^*=-J_A(u)\ne J_A(u^*).
\]
The first inequality follows already from the values at zero.
The second uses the reflection in the convolution involution.
Thus $J_A$ is neither an algebra homomorphism nor an involution-preserving
map.  Positivity of the numerical Gaussian average cannot pass through
it by Theorem~\ref{thm:dagger}.

\section{The remaining comparison with quantum currents}
\label{sec:weil-current-comparison}

The formal Hamiltonian bracket on $H_{\mathrm{for}}$ is
$[f,g]=f_xg_y-f_yg_x$ modulo constants.
Its continuous adjoint operators provide symmetry actions on formal
coefficient kernels.  A numerical positive representation needs further
data: a complex Hilbert space, a common invariant dense domain, and
operators whose restricted adjoints realize the chosen involution.
Proposition~\ref{prop:weil-formal-adjoint} prevents obtaining these data
by completing all finite matrix kernels in the formal operator topology.
Proposition~\ref{prop:weil-gaussian-comparison} also prevents specializing
an arbitrary formal parameter at a positive numerical value.
These are restrictions on the stated carriers and maps.

For a positive arithmetic interpretation, the first missing comparison
is an involution-preserving algebra map
\[
 T:\A_\zeta\longrightarrow\mathfrak B
\]
into a specified algebra of current observables, together with a positive
linear functional $\omega$ on $\mathfrak B$ satisfying
$L_\zeta=\omega\circ T$.
The algebra, involution, domain and positivity of $\omega$ must be
constructed independently of the zeta zero locations.
Theorem~\ref{thm:dagger} would give the isometry
$U:[f]\mapsto[T(f)]$ onto the closed subspace
$K=\overline{[T(\A_\zeta)]}\subset\mathcal H_\omega$.
Proposition~\ref{prop:graph-transfer} identifies the transported common
graph domain with that of the restricted multipliers on
$E=[T(\A_\zeta)]\subset K$, with every closure taken in $K$.
Equality with the full ambient common domain requires a further
hypothesis, such as surjectivity of $T$.
The finite-mass obstruction in Theorem~\ref{thm:weil-mass} still applies:
a finite-norm vector representing this functional, or a positive
functional on a $C^*$-algebra as target, cannot supply this factorization.
A positive functional on a proper nonunital algebra remains a possible
carrier, as the finite-matrix trace demonstrates.

Any current-field construction must also specify its operator products,
contraction tensors and the scalar trace or weight on those products.
Their convergence with the spatial propagators is an additional analytic
condition.  Neither the strict zero strip nor the equality of the
Gaussian moment formulas constructs this comparison or proves the
positivity of the full Weil form.
